
\section{The Claim under Test}

A familiar Christian reading joins three details. Jesus is born in Bethlehem,
whose name Jerome and Ambrose receive as ``house of bread.'' Mary wraps him in
swaddling cloths and lays him in a manger, a feeding place. Later he says, ``I
am the bread of life,'' and identifies the bread he will give as his flesh for
the life of the world (Jn 6:35, 51 in modern numbering; 6:35, 52 in the
Douay--Rheims/Vulgate numbering used by this repository). The conjunction is often presented as
though Luke and John had encoded one transparent Eucharistic sign.

The pieces are not equally established. Luke supplies the swaddling cloths and
the manger, but does not explain them as Eucharistic. John supplies bread from
heaven, flesh, eating, drinking, abiding, and resurrection, but does not mention
Bethlehem, the manger, or the cloths. The patristic ``house of bread'' gloss
does not prove that the name was coined as prophecy. The question is therefore not
whether the conjunction is theologically fruitful. It is: \emph{how early, how
directly, and with what limits did Christian readers make it?}

The answer presently supported by the checked record is layered. The
Eucharistic reading of John 6 is early and explicit across several major Greek
and Latin witnesses. A direct
Bethlehem--John 6 conjunction is securely visible in late-fourth-century Latin
reception. A manger-as-food image also occurs in ancient Christian preaching,
but one especially vivid text commonly assigned to the third century belongs
to a collection transmitted as doubtful or spurious. The swaddling cloths add
the Incarnation's poverty, visibility, and passibility. None of the exact
Nativity/reception witnesses named in this study turns the cloths themselves
into Eucharistic wrappings; that bounded result does not survey all patristic
literature.

